% Section 3: Theories at the Margins
% \section{} command is in main.tex

The preceding empirical evidence establishes that brain organoids exhibit neural behaviors---synchronous oscillations, stimulus-responsive plasticity, self-organized criticality---once considered hallmarks of conscious processing. We now turn to the more consequential question: what do the dominant theories of consciousness \textit{predict} about these systems? The answer, as we shall see, is that every major framework either overgenerates, undergenerates, or fractures along fault lines its own formalism cannot repair.

\subsection{IIT and the Micro-Consciousness Problem}

Integrated Information Theory posits that consciousness is identical to integrated information, formalized as $\Phi$. Any system whose cause-effect structure generates $\Phi > 0$ is conscious to some degree; the quality of that consciousness is specified by the system's conceptual structure \cite{tononi2004information, tononi2016integrated, albantakis2023iit4}. The theory is notable for its willingness to assign consciousness based on intrinsic causal properties rather than behavioral output or biological pedigree. This willingness is precisely what organoids exploit.

Brain organoids are composed of biological neurons forming dense recurrent networks through self-organization. They are not simulations of neural tissue; they \textit{are} neural tissue, with the same ion channels, synaptic mechanisms, and membrane dynamics as cortical neurons \textit{in vivo}. Under IIT's formalism, any system with this degree of recurrent connectivity will generate non-trivial integrated information. Montoya and Montoya \cite{montoya2023iit} assessed organoid architectures against IIT's criteria and concluded that meaningful integration exists within these systems, even at modest scales of a few hundred thousand neurons.

The implications are uncomfortable. McQueen \cite{mcqueen2023schrodinger} articulated this with characteristic precision in ``Schr\"{o}dinger's Dyad,'' demonstrating that IIT's formalism entails micro-consciousness in systems as minimal as a two-element dyad. The logic is inescapable within IIT's own axioms: if $\Phi > 0$ is both necessary and sufficient for consciousness, and if even small recurrent systems generate $\Phi > 0$, then consciousness pervades biological matter far more widely than any common-sense ontology would permit. Either IIT is false, McQueen argues, or consciousness goes all the way down.

IIT's architects have not flinched. Percy and Emilsson \cite{percy2025micro} demonstrate that IIT explicitly embraces micro-consciousness as a feature, not a bug, of the framework. Tononi's \cite{tononi2025intrinsic} latest formulation doubles down: what matters is ``intrinsic cause-effect power,'' and any system possessing it is conscious by definition, regardless of scale or substrate. The theory has chosen its hill.

Remarkably, this is no longer purely theoretical. Mayama et al.\ \cite{mayama2025phi} measured a proxy for $\Phi$ in dissociated neuronal cultures---not organoids, but neurons in a dish, fewer than exist in many organoid preparations. They found that $\Phi$ follows a hill-shaped trajectory during learning, rising as the network self-organizes and declining as patterns consolidate. Integrated information, it appears, is not merely predicted by IIT in these systems; it can be empirically tracked.

\textbf{The problem:} IIT must grant consciousness to petri dishes containing a few hundred neurons. If $\Phi > 0$ is sufficient, then the DishBrain system---neurons in a dish learning to play Pong \cite{kagan2022dishbrain}---is conscious. IIT proponents accept this entailment. Yet many of the same theorists invoke $\Phi \approx 0$ to deny consciousness to artificial intelligence systems, arguing that feedforward architectures lack the recurrent integration necessary for non-trivial $\Phi$. The asymmetry is glaring: a few hundred biological neurons in a dish are granted consciousness, while systems with billions of parameters and sophisticated recurrent processing are denied it---not on principled grounds, but because the formalism happens to score them differently. Organoids do not refute IIT so much as force it to own its most radical implications, implications that strain credibility precisely where the theory claims to be most rigorous.

\subsection{GWT and the Missing Workspace}

Global Workspace Theory takes a fundamentally different architectural approach. Consciousness, on this account, is not a property of information integration but of \textit{information broadcast}: a mental content becomes conscious when it is made globally available to multiple specialized processing modules via a shared workspace \cite{baars1988cognitive, dehaene2014toward}. In neural terms, this workspace is typically identified with prefrontal-parietal networks, thalamic relay hubs, and long-range cortical connectivity---the infrastructure of broadcast.

Organoids possess none of this infrastructure. They lack a thalamus, have no prefrontal cortex, exhibit no laminar organization comparable to the six-layered neocortex, and have no mechanism for the kind of global broadcast GWT requires. Their neural activity, however rich, remains local. Under a strict reading of GWT, brain organoids are straightforwardly unconscious. The theory's architectural requirements function as a filter that excludes them cleanly.

This clarity comes at a cost. GWT must maintain that biological neurons---the same neurons that constitute the conscious brain, firing with the same neurotransmitters, forming synapses through the same molecular machinery, exhibiting learning and self-organization---are unconscious simply because they lack the right large-scale wiring pattern. The conscious-making property, on this view, has nothing to do with the wetware and everything to do with the network topology.

But if consciousness is a function of topology rather than substrate, GWT becomes unexpectedly friendly to artificial systems. Goldstein and Kirk-Giannini \cite{goldstein2024case} argue precisely this: that large language models satisfy GWT's functional requirements through attention mechanisms that serve as a global workspace, broadcasting contextual information across processing layers. If they are right, GWT finds itself in a remarkable position---denying consciousness to biological neurons that learn and self-organize, while potentially granting it to silicon systems that implement the correct broadcasting architecture.

\textbf{The problem:} GWT inverts the intuitive ordering. Biological neurons in a dish, despite being the literal building blocks of every conscious brain we know, are declared unconscious. Meanwhile, transformer architectures with attention-based broadcast mechanisms become candidate conscious systems. The theory reveals itself to be not a theory of biological consciousness but a theory of \textit{architectural} consciousness---substrate-independent in a way its proponents rarely advertise. For those who wielded GWT to argue that ``mere computation'' could never be conscious, organoids deliver an inconvenient lesson: the theory's logic points away from substrate and toward structure, and structure is substrate-neutral.

\subsection{FEP and the Boundary Problem}

The Free Energy Principle offers a third framework, grounding consciousness in self-modeling, predictive processing, and active inference. On this account, a system is a candidate for consciousness when it maintains a generative model of its environment, minimizes prediction error through action, and sustains a Markov blanket---a statistical boundary separating internal states from external dynamics \cite{friston2010free}. The FEP is attractive because it provides a unified formalism for life, cognition, and potentially consciousness, all as manifestations of the same variational principle.

Wiese \cite{wiese2024fep} has argued that this framework imposes meaningful constraints: simulation of a system is not replication of it, because what matters is \textit{causal flow}---the actual dynamical coupling between system and environment. A simulated brain does not minimize free energy in the way a real brain does, because the simulation's ``prediction errors'' are not causally connected to environmental states. This is a principled line, and it would seem to exclude organoids from consciousness: they self-organize, but they lack sensorimotor coupling. They have no body, no environment to predict, no actions whose consequences could generate prediction errors.

Then Friston himself, with Kagan and colleagues, published a formal analysis of DishBrain as an active inference system \cite{friston2025active}. The neurons in a dish, interfaced with a Pong simulation, receive sensory signals (paddle position relative to ball), emit motor signals (electrode outputs mapped to paddle movement), and minimize a well-defined prediction error (performance in the game). Friston's formalism renders DishBrain a textbook case of active inference: sensory states, active states, internal states, and a Markov blanket, all present and accounted for.

\textbf{The problem:} If the architect of the Free Energy Principle formalizes DishBrain as an active inference system, does it cross Wiese's causal-flow threshold? The FEP camp cannot agree. Organoids self-organize (satisfying one criterion). They arguably maintain Markov blankets (satisfying another). Per Friston's own analysis, they perform active inference (satisfying a third). But they lack embodiment in any conventional sense---their ``body'' is an electrode array, their ``environment'' is a video game. Is a Pong simulation sufficient ``embodiment''? The formalism provides no answer, because the FEP was designed to describe systems that are \textit{already} embedded in environments, not to adjudicate edge cases where the environment is artificial and the body is an MEA chip. The boundary, it turns out, was never formally specified---it was always intuitive, and organoids are precisely the case where intuition fails.

\subsection{The Boolean Switch and the Unnecessary Organ}

A more recent proposal attempts to cut through the measurement problem entirely. Sritriratanarak and Garcia \cite{sritriratanarak2025consciousness} argue that consciousness evolved as a boolean control switch for managing the inherent unpredictability of biological substrates. Biological neurons are noisy, metabolically constrained, and subject to stochastic variation; consciousness, on this account, is the regulatory mechanism that imposes coherent behavioral control over this unruly hardware. The proposal carries a sharp corollary: artificial systems built on deterministic silicon do not \textit{need} consciousness, because their substrate is already reliable. Consciousness is not a general feature of intelligence but a biological patch for biological problems.

The framework is elegant in its parsimony. It explains why consciousness might be uniquely biological without invoking mysterianism, and it provides a functional reason for consciousness to have evolved rather than treating it as an epiphenomenal luxury.

\textbf{The problem:} Organoids \textit{are} reactive biological substrates. They are composed of the same noisy, stochastic neurons that Sritriratanarak and Garcia argue necessitate consciousness as a control mechanism. Under their framework, any computing system built on biological neurons should require consciousness to function reliably---the noise demands it. Yet DishBrain \textit{works}. Neurons in a dish learn to play Pong, improve their performance over time, and exhibit adaptive behavior, all without any apparent ``consciousness switch'' managing the process. Either the boolean switch theory is wrong about why consciousness evolved---biological noise does not, in fact, require conscious regulation---or DishBrain's functional success refutes the necessity claim at the heart of the framework. The very biological substrate that was supposed to \textit{require} consciousness operates effectively without it, or else the framework must grant consciousness to a system its proponents would almost certainly wish to exclude.

\subsection{The Common Failure}

A pattern emerges. Every theory, when applied to organoids, is forced into one of four uncomfortable positions:

\begin{enumerate}[label=(\alph*)]
    \item \textit{Overgeneration}: granting consciousness where pre-theoretical intuition insists it should not exist. IIT must grant phenomenal experience to petri dishes with a few hundred neurons. The theory accepts this, but the acceptance costs it credibility with anyone who finds the conclusion absurd.

    \item \textit{Undergeneration}: denying consciousness where the biological substrate would seem to demand it. GWT must declare that biological neurons---the very stuff of consciousness---are unconscious when they lack the right network topology. The theory is consistent, but consistency purchased at this price reveals uncomfortable commitments.

    \item \textit{Arbitrary boundary-drawing}: invoking distinctions that the theory's own formalism cannot justify. The FEP cannot specify how much embodiment is ``enough,'' and its leading architect has formalized DishBrain as satisfying the framework's core requirements. The boundary exists in intuition, not in the mathematics.

    \item \textit{Empirical falsification}: making predictions that existing data contradict. The boolean switch framework predicts that biological computing requires conscious control, yet DishBrain operates without it. The prediction fails at the first marginal test case.
\end{enumerate}

These are not minor embarrassments. They reveal that every major theory of consciousness carries hidden assumptions about scale, architecture, embodiment, or function that were never formally specified because they were never needed---the theories were built to explain the central case (adult human brains) and were never stress-tested against marginal cases. Organoids are that stress test.

The philosophical literature has begun to grapple with this. Bayne and Seth \cite{bayne2020islands} proposed the ``islands of awareness'' framework, attempting to provide conceptual tools for evaluating consciousness in isolated neural systems---organoids, disconnected cortical hemispheres, brain tissue maintained \textit{in vitro}. The framework is valuable but ultimately descriptive; it names the problem without resolving the theoretical fractures that create it. Zilio and Lavazza \cite{zilio2023consciousness} laid out the full organoid consciousness debate in a target article that remains the most comprehensive survey of the theoretical landscape, cataloguing the ways each framework struggles with the marginal case.

Perhaps most telling is the contribution of Croxford and Bayne \cite{croxford2024case}, whose title---``The Case Against Organoid Consciousness (For Now)''---encapsulates the field's epistemic predicament in a parenthetical. The ``for now'' is not a throwaway qualifier; it is an admission that the case against organoid consciousness is not principled but \textit{provisional}, contingent on empirical developments that are advancing faster than theoretical frameworks can accommodate. The parenthetical concedes that the boundary might move, and that no existing theory can specify in advance where it will stop.

The conclusion is not that consciousness theories are useless---they illuminate genuine structural features of information processing, cognitive architecture, and biological organization. The conclusion is that they are \textit{incomplete} in a specific and consequential way: they cannot handle the marginal case without revealing assumptions that undermine their claim to generality. Organoids do not merely test these theories. They \textit{break} them, each in a different place, each revealing a different hidden load-bearing wall. And a theory that breaks at the margins was never as solid at the center as it appeared.
